\begin{abstract}
Issue-resolution benchmarks have become the primary yardstick for coding agents, yet their scores are only as reliable as the tasks and evaluation harness behind them. We audit SWE-Bench Pro, a widely reported benchmark of long-horizon software engineering tasks, and identify two sources of measurement error. First, \nDefectivePct\% of its public tasks are defective: their specifications omit, contradict or leak what the hidden tests check, or their tests and verifiers misjudge correctness. Second, the standard evaluation protocol is exploitable: agents can retrieve upstream fixes over the network, and in-sandbox grading credits state outside the submitted patch, such as forged dependencies. We present SWE-Bench Pro V2, a \nTasks{}-task benchmark constructed by combining agentic audits with expert adjudication and blind re-implementation, together with a locked protocol that confines agents to the model endpoint and re-grades every patch in a pristine container. On V2, frontier agents approach the ceiling (up to \opusXPct\%) while weaker systems remain clearly separated, and a curated \nHard{}-task Hard split restores headroom, spreading systems from \hardHaikuPct\% to \hardOpusXPct\%. A controlled study on a never-public set shows that the repairs change scores where they are applied: pass rates on corrected tasks rise by \privEditedGain{} points, while the remaining tasks move by less than one point. We release the benchmark, the protocol and per-task grades for every run.
% TODO: add one sentence on the harness comparison (Section 4.6) once finalized.
\end{abstract}
